\section{Related Work}

\paragraph{Efficient diffusion inference.}
Diffusion models require repeated neural network evaluation over many denoising steps, and video generation further increases the cost through long spatial temporal token sequences. Existing acceleration methods reduce this cost through step reduction, quantization, token selection, feature caching, sparse attention, and adaptive computation. Recent approaches also coordinate multiple acceleration mechanisms to reduce computation across timesteps, tokens, and attention connections \citep{liu2025astraea,zhao2026rapid}. Most of these techniques optimize a single generation process and do not exploit reusable computation from previously completed requests.

\paragraph{Intra request acceleration.}
Intra-request methods exploit redundancy within one denoising trajectory. Feature caching reuses intermediate representations across nearby timesteps, while token wise methods selectively refresh tokens according to their sensitivity to caching errors \citep{zou2024accelerating,liu2025speca}. Sparse attention methods reduce quadratic attention cost by identifying structured patterns across timesteps, layers, and heads, then executing only selected query key interactions with efficient kernels \citep{chen2026sparse}. Diffusion attention has also been used to recover semantic regions from cross attention and self attention maps \citep{tian2024diffuse}. These methods provide important foundations for selective computation, but their reuse decisions are derived from states within the current trajectory. CARNATOR instead considers whether latent states, intermediate features, token regions, and attention patterns from a different generation trajectory remain compatible with the target request.

\paragraph{Inter request acceleration.}
Inter-request reuse has recently emerged as a new direction for diffusion serving. Early work retrieves similar text to image requests and reuses intermediate noise states to skip initial denoising computation \citep{agarwal2024approximate}. Subsequent methods extend this idea to text to video generation through entity level retrieval, shared semantic components, staged latent reuse, and region selective recomputation \citep{cherian2026chai,liu2026beyond}. Historical sparse attention patterns can also be transferred across related image to video requests to reduce online mask construction \citep{liu2026chorus}. 
CARNATOR formulates inter-request reuse as a generation-level compatibility problem.
It uses an early object-aware signature to filter historical caches according to the target generation state, and exploits accepted caches through joint region-restricted computation and historical sparse-connectivity reuse.
